% Conclusion and acknowledgments (Dr. Lin's structure: 0.5 page).
\section{Conclusion}
\label{sec:conclusion}

In this paper, we presented NOMAD, an in-network traffic obfuscation mechanism against ICS device fingerprinting over plain-text communication. NOMAD uses the syntax of control units to apply a different strategy to each type of operation on one programmable switch: it replaces the CLRT of READ and SELECT with a configured value, holds OPERATE commands inside the switch, and splits the control units of every response into one fixed layout. On a Tofino-1 testbed with 98,816 DNP3 transactions, the CLRT at the master stayed within 10--16~$\mu$s of the configured value at the median, and all responses carried 61~bytes and remained conformant to IEEE~1815. The CLRT classifier of Formby et al., trained before deployment, fell to chance in every setting, as it also did under a constant delay, and the same classifier retrained after deployment stayed near its unobfuscated accuracy on our software outstation, whose unobfuscated CLRT barely differs between operations. All transactions completed within 53~ms without retransmission. In future work, we will evaluate NOMAD with physical relays of several vendors, whose wider timing gives the classifier more to learn, and extend the strategies to requests.

\section*{Acknowledgments}
% To be completed by the authors.
